\section*{Data, summaries and distances}
\textbf{Shared notation:} $N$ = observations, $M$ = attributes; $i,\ell$ index observations, $j,m$ attributes, $c$ classes, $k$ components/clusters/folds as stated; $x_i$ = input row, $x_{ij}$ = its entry, $y_i$ = target, $\hat{\ }$ = estimated/predicted value.

\subsection*{Data matrix and standardization}
\textbf{Purpose:} Put attributes on comparable scales before distances, PCA or regularized models.
\[
X\in\R^{N\times M},\qquad \tilde x_{ij}=\frac{x_{ij}-\bar x_j}{s_j},
\]
\[
\bar x_j=\frac1N\sum_i x_{ij},\qquad
s_j^2=\frac1{N-1}\sum_i(x_{ij}-\bar x_j)^2 .
\]
\textbf{Symbols:} $X\in\R^{N\times M}$ is the data matrix; $\R$ means real values; $\bar x_j$ and $s_j$ are mean and sample standard deviation of attribute $j$; $\tilde X$ is standardized data.

\textbf{How to read it:} $\tilde x_{ij}>0$ means observation $i$ is above the mean of attribute $j$, measured in standard deviations.

\subsection*{Covariance, correlation and distance}
\textbf{Purpose:} Measure joint variation or similarity between observations.
\[
\hat\Sigma=\frac1{N-1}\tilde X^\top\tilde X,\qquad
\rho_{jk}=\frac{\Cov(x_j,x_k)}{s_js_k},
\]
\[
d(x_i,x_\ell)=\sqrt{\sum_m(x_{im}-x_{\ell m})^2}.
\]
\textbf{Symbols:} $\hat\Sigma$ = sample covariance/correlation matrix of standardized attributes; $\rho_{jk}$ = correlation of columns $j,k$; $\Cov$ = covariance; $d(x_i,x_\ell)$ = Euclidean distance between rows.
\textbf{Means:} covariance/correlation compare columns; distance compares rows. For standardized $X$, the covariance matrix equals the correlation matrix.

\textbf{How to read a matrix:} diagonal correlation is $1$; sign gives direction of linear relation and $|\rho|$ its strength. Distance matrices have zero diagonal and large entries for separated observations.

\[
d_1(x,z)=\sum_m|x_m-z_m|,\qquad
d_\infty(x,z)=\max_m|x_m-z_m|.
\]
\textbf{Symbols/use:} $x,z$ = two vectors; $d_1$ is cityblock distance; $d_\infty$ is maximum-coordinate distance. Use the named distance in neighbour/density calculations.

\subsection*{Plots and summary objects}
\begin{tabularx}{\linewidth}{@{}lX@{}}
\toprule
\textbf{Object} & \textbf{How to read it}\\ \midrule
Histogram/density & Height shows where values concentrate; width/tails show spread and extremes.\\
Boxplot & Middle line is median; box is middle $50\%$; whiskers/points show range or unusual values.\\
Scatter plot & Tilt indicates correlation sign; narrow cloud indicates strong linear relation.\\
Scatter matrix & Panel $(j,k)$ compares two attributes; match signs to a proposed correlation matrix.\\
\bottomrule
\end{tabularx}
\textbf{Do this:} identify variables and scale; read direction/spread; then match to numerical summaries.

\subsection*{Regression loss}
\[
e_i=y_i-\hat y_i,\quad
\text{SSE}=\sum_i e_i^2,\quad
\text{MSE}=\frac1N\sum_i e_i^2,
\]
\[
\text{RMSE}=\sqrt{\text{MSE}}.
\]
\textbf{Symbols/means:} $e_i$ = residual for case $i$; $y_i,\hat y_i$ = true and predicted response; SSE = summed squared error; MSE = its mean; RMSE = its square root in the unit of $y$. Smaller test/CV RMSE indicates better prediction.
